\section{Discussion}
\label{sec:discussion}

\paragraph{How our predictions fared.} Of the predictions in \S\ref{sec:method}, RQ2
and RQ4 held as written. RQ1 held in its shape but not in its size: Phase~1 costs
68--79\,\textmu s per trace line, not 10--30, because of the one-byte reads we had not
anticipated; Phase~2's per-file cost (7.2\,ms) and its growth with the tree held. RQ3
held on ext4 and XFS and failed on F2FS, whose lookups cost 0.4\,ms even warm and need
one device read per file even for files created together. The interpreter-lock result
of \S\ref{sec:gil} was not predicted: we found the cause while explaining why threads
never helped in E4, and designed E8 afterwards to test it, so it should be read as
exploratory, as should everything in \S\ref{sec:explore}.

\paragraph{When could threads pay?}
Threads overlap waiting, so they help only when the analysis waits on many
independent device reads that the device can serve concurrently. In FragPicker's
workflow that is the exception: the analysis runs right after tracing the
application, which has just touched the same inodes, directories and extent blocks,
so the lookups hit the cache (\S\ref{sec:fiemap}). The cases that remain are an
analysis run long after the trace, on a machine whose memory cannot hold the
metadata, over files whose inodes are scattered, on a device with deep queues
(network block storage, a RAID of disks). There the waits are in the tree walk that
builds the inode map, which the fork runs as one \texttt{find} before any parallel
work (\S\ref{sec:decompose}), and a walk split across a few threads does overlap them:
3.9$\times$ under dm-delay and 6.1$\times$ on a real network disk (\S\ref{sec:walk},
\S\ref{sec:realdisk}). On CPython 3.14 the per-file lookups would overlap as well
(\S\ref{sec:gilpy}), but in the common case they have nothing to wait for.

\paragraph{What the analysis phase should do instead.}
Our measurements point to three changes, in order of benefit, none of which needs
parallelism.
\begin{enumerate}
\item \emph{Do Phase~1 in memory.} Grouping requests by inode in a dictionary and
  writing each request file once removes three system calls per trace line; this is
  most of the rewrite's gain on long traces.
\item \emph{Never resolve inodes one \texttt{find} at a time.} One walk of the tree
  replaces $F$ walks. Better still, the tracer could record the device number and
  the inode generation next to the inode number; then a file can be opened with
  \texttt{open\_by\_handle\_at} without any path lookup, and inodes from other file
  systems can no longer be confused with files under \texttt{/mnt}.
\item \emph{Look extents up once.} Extents are needed only by the migration phase,
  which already reads them; reading them during analysis as well, as the parallel
  path does, doubles the work.
\end{enumerate}
Aggregating requests in a BPF map instead of printing every event would also remove
the two tracing problems we found, the lost tail of the output and the dropped
samples, and the 32-bit position format; we leave that design to future work.

\paragraph{Lessons for evaluating parallel tools.}
The 12.74$\times$ result was produced by the latency model, not by the code. A
\texttt{sleep} before every lookup charges the same latency to cached and uncached
lookups and lets all of them wait at once. Three habits would have caught it: compare
against the best sequential implementation, not against the parallel engine with one
thread; inject latency below the cache, at the block layer, so that only real device
reads are slowed; and count device reads per unit of work, which tells in advance how
much there is to overlap.
